% Piezas que comparten los tres temas nuevos. La carga estilo.latex justo antes
% de \begin{document}, después de paquetes.tex, estilos.tex y de lo que pandoc
% emite, así que aquí se redefine lo que ya existe.
\usepackage{eso-pic,pdfrender,anyfontsize}
\tcbuselibrary{skins,breakable}
\usetikzlibrary{arrows.meta}
% Los entornos de amsthm de estilos.tex y comandos.tex se envuelven en una caja;
% numeración y cabecera siguen siendo las suyas.
\newcommand\temaenvolver[2]{\foreach \e in {#1}{\expandafter\tcolorboxenvironment\expandafter{\e}{#2}}}
% Título partido: la primera palabra en macizo y el resto aparte (\temaresto).
\def\temapartir#1 #2\relax{\temaprimera{#1}\ifblank{#2}{}{\temaresto{#2}}}
\providecommand\temaprimera[1]{#1}\providecommand\temaresto[1]{ #1}
